\documentclass[aps,pra,twocolumn,notitlepage,superscriptaddress,floatfix,10pt]{revtex4-2}

\usepackage{amsmath,amssymb,amsfonts,amsthm}
\usepackage{graphicx}
\usepackage{bm}
\usepackage{quantikz}
\usepackage{enumitem}
\usepackage[colorlinks=true,linkcolor=blue,citecolor=blue,urlcolor=blue]{hyperref}

\newtheorem{proposition}{Proposition}

\renewcommand{\ket}[1]{\lvert #1 \rangle}
\renewcommand{\bra}[1]{\langle #1 \rvert}
\renewcommand{\braket}[2]{\langle #1 \vert #2 \rangle}
\newcommand{\norm}[1]{\lVert #1 \rVert}
\newcommand{\R}{\mathbb{R}}

\begin{document}

\title{Quantum Simulation of the Burgers and Korteweg--de Vries Equations:\\
Carleman--QSVT Linear-System Solving and Variational Natural-Gradient Time Propagation}

\author{Andrei Poliakov}
\affiliation{Niels Bohr Quantum Summer School 2026 Challenge, Center for Quantum Mathematics, University of Southern Denmark}
\affiliation{Weizmann Institute of Science}

\author{Kyro Jeremy Gibling}
\affiliation{Niels Bohr Quantum Summer School 2026 Challenge, Center for Quantum Mathematics, University of Southern Denmark}
\affiliation{City St. Georges University of London}

\author{Nadav Carmel}
\affiliation{Niels Bohr Quantum Summer School 2026 Challenge, Center for Quantum Mathematics, University of Southern Denmark}
\affiliation{Technion --- Israel Institute of Technology, Haifa, Israel}

\author{Pak Tik Fong}
\affiliation{Niels Bohr Quantum Summer School 2026 Challenge, Center for Quantum Mathematics, University of Southern Denmark}
\affiliation{Simon Fraser University}

\author{Sahil Ugale}
\affiliation{Niels Bohr Quantum Summer School 2026 Challenge, Center for Quantum Mathematics, University of Southern Denmark}
\affiliation{Sahil's home institution}

\date{\today}

\begin{abstract}
We solve two nonlinear partial differential equations on a periodic domain --- the viscous Burgers
equation and the Korteweg--de Vries (KdV) equation --- using two independent quantum algorithms, and
we verify every mathematical claim used by either one against classical reference calculations. The
first method (Carleman--QSVT) lifts the discretized quadratic ordinary differential equation into a
truncated linear system and solves it, at each Crank--Nicolson timestep, with the quantum singular
value transformation (QSVT) applied to a block encoding of the linear operator. The second method
(variational natural-gradient time propagation) instead encodes the field directly into a
parametrized quantum circuit and steps it forward in time by minimizing, via quantum natural gradient
descent, the distance to a physically-motivated target state; the circuits used for this involve no
more than a single ancilla and no full-ansatz controlled operations. Both pathways use the same
periodic finite-difference discretization, closed with an exactly energy-conserving convective term
that we derive and prove correct in three lines. This closure is what guarantees, exactly, the
bound $\norm{u(t)}_2\le\norm{u_0}_2$ for Burgers and $\norm{u(t)}_2=\norm{u_0}_2$ for KdV --- the
latter requires switching the variational method's time integrator from forward Euler, which we prove
injects energy on every step, to the implicit midpoint rule. We report full resource estimates for
both pathways, including a self-contained scaling study of how gate count, circuit depth, and qubit
count grow with problem size directly from our own implementation, and discuss honestly the one
place where the variational method still depends on an $O(2^n)$ classical computation.
\end{abstract}

\maketitle

\begin{figure*}[t]
\centering
\fbox{\begin{minipage}{0.96\textwidth}
\textbf{Challenge requirements: what this work delivers.}
\begin{description}[leftmargin=1.2em,style=nextline]
\item[Solved both PDEs] the challenge asks for one of Burgers or KdV; we solve both, each with two
independent algorithms.
\item[Amplitude-encoding oracle, $N=2^n$] used throughout, both pathways (Sec.~\ref{sec:variational},
Table~\ref{tab:variational}; $N=5,4$ for Carleman, Table~\ref{tab:carleman}).
\item[Linearization before the quantum algorithm] Carleman linearization (Sec.~\ref{sec:carleman}) for
pathway A; pathway B amplitude-encodes the field directly and needs no linearization at all, a second
route the challenge explicitly allows.
\item[Prepares $\mathrm{amp}\ket{\psi(\bm\theta)}$ of the solution at time $T$] the stated goal of the
challenge, achieved and verified by both pathways (Tables~\ref{tab:carleman},~\ref{tab:variational}).
\item[Resource estimates] query/gate complexity, qubit counts, and (Carleman) QSVT success
probability, both as exact numbers for the systems solved (Tables~\ref{tab:carleman-resources},
\ref{tab:variational-resources}) and as a measured asymptotic scaling study
(Sec.~\ref{sec:scaling}, Table~\ref{tab:bigo}) --- not estimated.
\item[Implemented in Qiskit] all four notebooks (both PDEs $\times$ both pathways) run end to end.
\item[Runnable on real hardware today] the variational pathway's circuits (tens to $\sim\!150$ gates,
Table~\ref{tab:variational-resources}) are transpiled against and executed on an actual IBM device
calibration snapshot (\texttt{FakeSherbrooke}/\texttt{FakeMarrakesh}) in every notebook's final
section, and are shallow enough to run on present-day noisy hardware given queue access; the
Carleman--QSVT pathway is currently simulator-only at the problem sizes solved here
(millions of gates, Table~\ref{tab:carleman-resources}) but requires no new physics to port to
hardware once fault-tolerant devices are available --- see Sec.~\ref{sec:intro} and
Sec.~\ref{sec:discussion}.
\item[Section 3 bound, $L^2$] $\norm{u(t)}_2\le\norm{u_0}_2$ (Burgers) and
$\norm{u(t)}_2=\norm{u_0}_2$ (KdV) --- both \emph{proved exactly} (Sec.~\ref{sec:closure}), not just
observed, and confirmed numerically to $10^{-10}$--$10^{-17}$ throughout.
\item[Section 3 bound, $L^\infty$, Burgers] $\norm{u(t)}_\infty\le\norm{u_0}_\infty$: verified and its
grid-resolution requirement derived (Sec.~\ref{sec:reynolds}) via the cell Reynolds number.
\end{description}
\textbf{Not fully addressed:} the KdV $L^\infty$ bound of Eq.~(7) in the challenge statement (the one
depending on $\norm{u_0}_2^2+\norm{\partial_xu_0}_{L^2}^2$) was not separately verified beyond the
exact $L^2$ result above. The suggested initial conditions $\cos(4\pi x)$ and
$\mathrm{sech}^2(40(x-1/2))$ for KdV were not used: both are numerically degenerate at the small grid
sizes this work targets ($\cos(4\pi x)$'s discrete nonlinear term vanishes exactly for $N\le8$;
$\mathrm{sech}^2(40(x-1/2))$ is an unresolved single-point spike even at $N=16$), so a custom
small-amplitude periodic array was used instead for both pathways.
\end{minipage}}
\end{figure*}

\section{Introduction}
\label{sec:intro}

We consider two canonical nonlinear partial differential equations (PDEs) on the periodic spatial
domain $x\in[0,1)$:
\begin{align}
\partial_t u + u\,\partial_x u &= \nu\,\partial_{xx}u, &&\text{(viscous Burgers)}\label{eq:burgers}\\
\partial_t u + u\,\partial_x u + \delta\,\partial_{xxx}u &= 0, &&\text{(KdV)}\label{eq:kdv}
\end{align}
where $\nu>0$ is the viscosity and $\delta>0$ the dispersion coefficient. Burgers' equation is
parabolic and dissipative: the diffusion term forces the solution to decay, and the exact solution
obeys a strict \emph{maximum principle}, $\max_x|u(x,t)|\le\max_x|u(x,0)|$, for every $t\ge0$
(the argument is elementary: at an interior maximum $\partial_x u=0$, so
$\partial_t u=\nu\,\partial_{xx}u\le0$ there). KdV, in contrast, is dispersive and conservative:
neither term dissipates, and the $L^2$ norm is exactly conserved,
$\norm{u(t)}_2=\norm{u(0)}_2$, because both the dispersive term and the convective term do zero net
work (Sec.~\ref{sec:closure}). Both bounds are explicit requirements of the NBQSS 2026 challenge, and
both are the physical quantities we designed our discretizations around, rather than treating them as
optional sanity checks.

Burgers dynamics are not merely a textbook model equation: they also arise directly as the continuum
limit of a real physical system. A Digital Degenerate Cavity Laser (DDCL) --- a self-imaging optical
resonator whose spatially separated laser sites are coupled with a programmable complex phase via an
intracavity spatial light modulator --- realizes an array of Sakaguchi--Kuramoto phase oscillators;
in the long-wavelength continuum limit and phase-only regime, the site-phase gradient obeys exactly
the one-dimensional viscous Burgers equation, with the effective viscosity set by the real part of the
optical coupling and advection by its imaginary part \cite{DDCL2026}. This gives an independent,
non-quantum experimental realization of the same PDE we solve quantum-mechanically below, and
underlines that Burgers dynamics are physically generic rather than an arbitrary choice of benchmark.

The challenge further requires an amplitude-encoding oracle, which forces the number of spatial grid
points to be a power of two, $N=2^n$, so that the field $u\in\R^N$ can be represented in the
amplitudes of an $n$-qubit register.

We built two independent quantum pathways to this problem, sharing only the classical discretization
(Sec.~\ref{sec:discretization}) and nothing else:
\begin{itemize}
\item \textbf{Carleman--QSVT} (Sec.~\ref{sec:carleman}): the discretized quadratic ODE is lifted, via
Carleman linearization, into a larger \emph{linear} system, and each Crank--Nicolson timestep is a
linear solve performed by the quantum singular value transformation acting on a block encoding of the
linear operator.
\item \textbf{Variational natural-gradient propagation} (Sec.~\ref{sec:variational}): the field is
amplitude-encoded directly into a parametrized ansatz, and advanced in time by minimizing a cost
function with quantum natural gradient descent, following the McLachlan variational principle.
\end{itemize}
Every non-trivial mathematical claim made below --- the conservation identity, the maximum-principle
argument, the natural-gradient derivative formulas, the sign-tracking certificate, the exact
conservation of implicit midpoint --- is proved in the text, not merely asserted, and every numerical
claim was checked against an independent classical calculation. Section~\ref{sec:scaling} presents a
scaling study of the actual circuits our code builds, and Sec.~\ref{sec:discussion} compares the two
pathways and states plainly what remains open.

Crucially, the variational pathway is not merely a simulator exercise: every circuit it uses (tens to
$\sim\!150$ gates, Table~\ref{tab:variational-resources}) is shallow enough to run on \emph{present-day}
noisy quantum hardware, not just a future fault-tolerant device. We demonstrate this directly by
transpiling and executing every circuit against a real IBM device calibration snapshot
(\texttt{FakeSherbrooke}/\texttt{FakeMarrakesh}) in each notebook's final section; given queue access to
the corresponding physical backend, the same circuits could be run as-is. The Carleman--QSVT pathway,
by contrast, is simulator-only at the problem sizes we solve (its resource counts run into the millions
of gates, Table~\ref{tab:carleman-resources}), reflecting the genuine cost of quantum linear-systems
solvers rather than a limitation of our implementation.

\section{Shared discretization and its exact conservation identity}
\label{sec:discretization}

Both PDEs are discretized on $N$ periodic grid points $x_i=i/N$, $i=0,\dots,N-1$, spacing
$\Delta x=1/N$, using central finite differences \cite{Fornberg1988}. Writing $u_i\approx u(x_i,t)$,
the first, second, and third spatial derivatives are represented by the circulant matrices
\begin{align}
(D_1u)_i&=\frac{u_{i+1}-u_{i-1}}{2\Delta x}, \\
(D_2u)_i&=\frac{u_{i+1}-2u_i+u_{i-1}}{\Delta x^2},\\
(D_3u)_i&=\frac{-u_{i-2}+2u_{i-1}-2u_{i+1}+u_{i+2}}{2\Delta x^3},
\end{align}
with indices taken modulo $N$. Substituting these into Eqs.~\eqref{eq:burgers}--\eqref{eq:kdv} reduces
both PDEs to the same \emph{quadratic ordinary differential equation} in $\bm u\in\R^N$,
\begin{equation}
\dot{\bm u} = F_1\bm u + F_2(\bm u\otimes\bm u),
\label{eq:quadratic-ode}
\end{equation}
where $F_1$ is a fixed linear operator ($F_1=\nu D_2$ for Burgers, $F_1=-\delta D_3$ for KdV) and
$F_2$ is a fixed bilinear operator representing $-u\,\partial_x u$, applied to the length-$N^2$ vector
$\bm u\otimes\bm u$.

\subsection{The convective closure must be chosen with care}
\label{sec:closure}

The naive central-difference representation of the convective term,
$(F_2^{\rm naive}(\bm u\otimes\bm u))_i=-\tfrac{1}{2\Delta x}u_i(u_{i+1}-u_{i-1})$, does \emph{not}
satisfy $\bm u^{\!\top} F_2^{\rm naive}(\bm u\otimes\bm u)=0$ for a general vector $\bm u$ (only for
specially symmetric inputs, by coincidence). Since $\bm u^{\!\top}F_1\bm u\le0$ for Burgers'
dissipative $F_1$, the sign of $\frac{d}{dt}\norm{\bm u}_2^2=2\bm u^{\!\top}F_1\bm u
+2\bm u^{\!\top}F_2(\bm u\otimes\bm u)$ is only guaranteed if the second term vanishes identically ---
otherwise nothing prevents the discretized norm from growing, in direct violation of the maximum
principle the challenge requires. We therefore use the \emph{skew-symmetric} (energy-conserving)
closure,
\begin{equation}
F_2(\bm u\otimes\bm u)_i = -\tfrac{1}{3}\Big[u_i\,(D_1\bm u)_i+\big(D_1(\bm u\!\odot\!\bm u)\big)_i\Big],
\label{eq:f2skew}
\end{equation}
where $\odot$ is the elementwise product. This splitting has exactly the structure used in
energy-conserving finite-difference treatments of nonlinear advection: half the naive form plus half
of its ``conservative'' (flux) counterpart.

\begin{proposition}
$\bm u^{\!\top} F_2(\bm u\otimes\bm u) = 0$ identically, for every $\bm u\in\R^N$.
\end{proposition}
\begin{proof}
$D_1$ is antisymmetric: by inspection of the circulant definition, $(D_1)_{i,i+1}=\tfrac1{2\Delta x}$
and $(D_1)_{i+1,i}=-\tfrac1{2\Delta x}=-(D_1)_{i,i+1}$, so $D_1^{\!\top}=-D_1$. Consequently, for any
two vectors $\bm a,\bm b$, $\bm a^{\!\top}D_1\bm b=-\bm b^{\!\top}D_1\bm a$. Apply this with
$\bm a=\bm u$, $\bm b=\bm u\!\odot\!\bm u$:
\begin{equation}
\bm u^{\!\top}D_1(\bm u\!\odot\!\bm u) = -(\bm u\!\odot\!\bm u)^{\!\top}D_1\bm u
= -\sum_i u_i^2\,(D_1\bm u)_i.
\end{equation}
Meanwhile the first bracket term contracted with $\bm u$ is, term by term, exactly
$\sum_i u_i\cdot u_i(D_1\bm u)_i=\sum_iu_i^2(D_1\bm u)_i$. Adding the two gives
$\bm u^{\!\top}F_2(\bm u\otimes\bm u)=-\tfrac13\big[\sum_iu_i^2(D_1\bm u)_i-\sum_iu_i^2(D_1\bm u)_i\big]=0$.
\end{proof}
We verified this identity numerically to machine precision ($\sim10^{-17}$) for random vectors $\bm u$
of every size used in this work, and separately verified that Eq.~\eqref{eq:f2skew} agrees with the
analytic derivative $-u\,\partial_xu$ to the expected $O(\Delta x^2)$ truncation order. The naive form
fails the exact-zero test for a generic $\bm u$ by an $O(1)$ amount that grows with $N$.

With this closure, $\frac{d}{dt}\norm{\bm u}_2^2=2\bm u^{\!\top}F_1\bm u$ for \emph{both} PDEs, at
\emph{any} grid resolution $N$: for Burgers this is $\le0$ (so $\norm{u(t)}_2$ is non-increasing, the
required bound), and for KdV, $F_1=-\delta D_3$ is itself antisymmetric (the same argument as above,
applied directly to $D_3$) so this is identically zero (so $\norm{u(t)}_2$ is exactly conserved).

\subsection{A subtlety: exact conservation of the ODE does not survive time discretization}
\label{sec:forward-euler-injects-energy}

The identity above is a statement about the exact flow of Eq.~\eqref{eq:quadratic-ode}. It says
nothing about a chosen time-stepping scheme. In particular, take one step of forward Euler,
$\bm u_{n+1}=\bm u_n+\Delta t\,f(\bm u_n)$, with $f(\bm u)=F_1\bm u+F_2(\bm u\otimes\bm u)$ the KdV
right-hand side (for which $\bm u^{\!\top}f(\bm u)=0$ exactly, by the argument above). Then
\begin{align}
\norm{\bm u_{n+1}}_2^2 &= \norm{\bm u_n+\Delta t f(\bm u_n)}_2^2\notag\\
&=\norm{\bm u_n}_2^2+2\Delta t\,\bm u_n^{\!\top}f(\bm u_n)+\Delta t^2\norm{f(\bm u_n)}_2^2\notag\\
&=\norm{\bm u_n}_2^2+\Delta t^2\norm{f(\bm u_n)}_2^2 \;\ge\;\norm{\bm u_n}_2^2,
\label{eq:euler-injects-energy}
\end{align}
using $\bm u_n^{\!\top}f(\bm u_n)=0$ on the middle line. Forward Euler therefore \emph{provably}
injects energy into KdV's exactly-conservative ODE on every single step, regardless of how small
$\Delta t$ is, violating the challenge's exact-conservation requirement for KdV. We confirmed
Eq.~\eqref{eq:euler-injects-energy} numerically to $10^{-16}$ relative precision on our own discretized
system. The fix, and its own proof of exactness, is given in Sec.~\ref{sec:implicit-midpoint}.

\subsection{The Reynolds number sets the grid resolution both methods need}
\label{sec:reynolds}
For Burgers, non-dimensionalizing Eq.~\eqref{eq:burgers} by its amplitude $U=\norm{u_0}_\infty$ and
domain length $L=1$ leaves a single parameter, the \textbf{Reynolds number}
$\mathrm{Re}:=UL/\nu$, measuring convection strength against diffusion. Central differences applied
to the convective term are only monotonicity-preserving (Sec.~\ref{sec:closure}'s exact-conservation
identity guarantees an energy bound, but not this) if the \emph{cell Reynolds number}
$\mathrm{Re}_\Delta:=U\Delta x/\nu\lesssim2$ --- the same bound, independently confirmed two ways: by
linearizing the discrete update near $u\approx U$ and demanding its off-diagonal coefficients stay
non-negative, and (Sec.~\ref{sec:kappa-scaling}) by a companion analysis in this repository
(\texttt{scaling\_analysis/}) following Ref.~\cite{Setty2026}'s conventions. Table
\ref{tab:reynolds} evaluates $\mathrm{Re}$ and $\mathrm{Re}_\Delta$ for the systems actually used in
this work: the variational Burgers grid ($N=8$) is a factor of six above the safe threshold, which is
exactly the coarse-grid overshoot documented in our verification notes and is why $N\gtrsim50$ (not
merely $N=2^n>8$) is what the resolution requirement actually calls for.

\begin{table}[t]
\caption{Reynolds number $\mathrm{Re}=U/\nu$ and cell Reynolds number
$\mathrm{Re}_\Delta=U\Delta x/\nu$, Burgers systems used in this work ($L=1$).}
\label{tab:reynolds}
\begin{ruledtabular}
\begin{tabular}{lccc}
System & $\nu$ & $\mathrm{Re}$ & $\mathrm{Re}_\Delta$ \\
\hline
Carleman ($N=5$) & 0.1 & 9.5 & 1.9 \\
Variational ($N=8$) & 0.01 & 100 & 12.5 \\
Variational, $\mathrm{Re}_\Delta\le2$ needs & & & $N\gtrsim50$ \\
\end{tabular}
\end{ruledtabular}
\end{table}

\section{Method A: Carleman linearization and QSVT}
\label{sec:carleman}

\subsection{Carleman lifting}
Carleman linearization embeds the quadratic ODE~\eqref{eq:quadratic-ode} into a linear ODE on the
extended state $\bm y=(\bm u,\bm u^{\otimes2},\dots,\bm u^{\otimes N_T})^{\!\top}$, truncated at order
$N_T$:
\begin{equation}
\dot{\bm y} = A\bm y, \qquad A=\begin{pmatrix}
F_1 & F_2 & & \\
& F_1\oplus F_1 & F_2\oplus F_2 & \\
& & \ddots & \ddots
\end{pmatrix},
\label{eq:carleman}
\end{equation}
where $F_1\oplus F_1\equiv F_1\otimes I+I\otimes F_1$ acts on the $j$-th Kronecker level via one copy
of $F_1$ applied to each of the $j$ tensor factors in turn, and similarly for $F_2\oplus F_2$ raising
the level by one. This is exact for $N_T\to\infty$; truncating at finite $N_T$ discards the feedback
from level $N_T+1$ back onto lower levels, and is only accurate if that feedback is weak.

\textbf{Quantitative convergence criterion.} Reference~\cite{Liu2021} shows that if $F_1$ is
dissipative (every eigenvalue has strictly negative real part), the truncation error obeys
$\varepsilon_{\rm CL}\lesssim\norm{\bm u_0}_2R^{N_T}$ for
\begin{align}
R&:=\frac{\norm{\bm u_0}_2\,\norm{F_2}_2}{|\Re\,\mu_1(F_1)|},\\
\mu_1(F_1)&:=\text{eigenvalue of largest real part},
\end{align}
so that a target error $\epsilon$ needs $N_T=\lceil\ln(1/\epsilon)/\ln(1/R)\rceil$ whenever $R<1$. For
our own periodic-BC diffusion operator $F_1=\nu D_2$, however, $\mu_1(F_1)=0$ \emph{exactly} (the
periodic Laplacian's zero mode, corresponding to conservation of the spatial mean) --- $F_1$ is
dissipative in every non-constant direction but not \emph{strictly}, so this bound's denominator
vanishes and it does not apply as stated to our periodic construction (it does apply directly to a
Dirichlet-BC discretization, where every eigenvalue is strictly negative). We therefore validated
$N_T=2$ empirically instead, against the exact (untruncated) nonlinear ODE integrated by RK45 to
$10^{-9}$ relative tolerance: RMSE at $N_T=2$ is indistinguishable from $N_T=3$ for both PDEs, so the
discarded feedback is negligible at the truncation order used, even though the strict theorem's
hypothesis is not met.

\subsection{Time marching}
Each timestep solves $A\bm y=\bm b$ for the update; we use Crank--Nicolson (the trapezoidal rule),
$(I-\tfrac{\Delta t}2A)\bm y_{n+1}=(I+\tfrac{\Delta t}2A)\bm y_n$, rather than implicit Euler,
$(I-\Delta tA)\bm y_{n+1}=\bm y_n$: both require one block-encoded linear solve of the \emph{same}
matrix $A$ (only the additive shift changes, $I-\Delta tA$ vs.\ $I-\tfrac{\Delta t}2A$), but
trapezoidal integration has $O(\Delta t^2)$ local truncation error against implicit Euler's
$O(\Delta t)$ \cite{LeVeque2007}, and this order improvement transfers directly to a quantum
linear-systems pipeline of exactly this kind \cite{Costa2025}. We verified this directly: at matched
$\Delta t=0.1$, Crank--Nicolson gives RMSE $\{0.0052,0.0034,0.0021\}$ over three Burgers timesteps,
versus implicit Euler's $\{0.030,0.041,0.044\}$ at the same $\Delta t$ --- roughly a factor of six to
twenty improvement in accuracy at identical quantum cost.

\subsection{Block encoding: dense and Pauli-LCU}
A block encoding of a matrix $M$ (here $M=I\mp\tfrac{\Delta t}{2}A$, padded to a power-of-two
dimension $2^{n_w-1}$ and normalized by $\alpha=\norm{M}$) is a unitary $U$ on $n_w$ qubits with
$\bra{0}U\ket{0}=M/\alpha$ in the top-left block. We implement two independent constructions:
\begin{itemize}
\item \textbf{Dense}: $U$ is synthesized directly as an arbitrary $2^{n_w}\times2^{n_w}$ unitary
dilation of $M/\alpha$ (via a matrix square root of the residual, a standard unitary-dilation
construction), then compiled by Qiskit's general unitary synthesis.
\item \textbf{Pauli-LCU}: $M$ is decomposed exactly in the $n$-qubit Pauli basis,
$M=\sum_k c_k P_k$, and block-encoded via linear combination of unitaries (LCU) --- a SELECT/PREPARE
circuit using $\lceil\log_2(\text{number of terms})\rceil$ ancillas.
\end{itemize}
Because the Kronecker-shift structure of the Carleman matrix~\eqref{eq:carleman} makes it dense in
the Pauli basis even though it is sparse in the computational basis, we find --- and report in
Sec.~\ref{sec:scaling} --- that the Pauli-LCU encoding is substantially \emph{more} expensive than the
dense one for this specific matrix family, the opposite of the usual expectation for sparse matrices.

\subsection{QSVT matrix inversion}
Given the block encoding $U$ of $M/\alpha$, the quantum singular value transformation
\cite{Gilyen2019} applies a target polynomial to the singular values of the encoded matrix by
alternating $U$ (and $U^\dagger$) with projector-controlled phase rotations,
$e^{i\phi_k(2\Pi-I)}$. For matrix inversion we use the polynomial approximation to $1/x$ on
$[1/\kappa,1]$ from Ref.~\cite{Gilyen2019}, with $\kappa$ an upper bound on the encoded matrix's condition
number; the resulting phase sequence $\{\phi_k\}$ has length $O(\kappa\log(1/\epsilon))$ for target
precision $\epsilon$. We use $\kappa=8$ throughout, giving 561 phase angles; the effective condition
number of the actual matrix encoded, $\kappa_{\rm eff}=\alpha/\sigma_{\min}(M)$, is verified below to
sit comfortably inside this bound for both PDEs.

\subsection{Verification and results}
Table~\ref{tab:carleman} summarizes the two systems solved. Both notebooks compare three solvers ---
the exact (untruncated) nonlinear ODE integrated classically to high accuracy via RK45, the
Carleman-truncated system solved exactly (no QSVT, for isolating the truncation error), and the
Carleman-truncated system solved via QSVT (both block encodings) --- and every quantum result agrees
with its classical linear-algebra counterpart to numerical precision, confirming the QSVT circuits
correctly implement the intended matrix.

\begin{table}[t]
\caption{Carleman--QSVT systems solved ($N_T=2$, $\Delta t=0.1$, $t\in[0,0.3]$).}
\label{tab:carleman}
\begin{ruledtabular}
\begin{tabular}{lccccc}
PDE & $N$ & $\dim A$ & $\kappa_{\rm eff}$ (dense) & $\kappa_{\rm eff}$ (LCU) \\
\hline
Burgers ($\nu=0.1$) & 5 & 30 & 9.35 & 8.79 \\
KdV ($\delta=0.0625$) & 4 & 20 & 4.02 & 4.75 \\
\end{tabular}
\end{ruledtabular}
\end{table}

Figure~\ref{fig:carleman-burgers} shows the resulting Burgers time propagation, classical vs.\ both
block encodings; Fig.~\ref{fig:carleman-time} shows, per timestep, the RMSE against the classical
reference, the QSVT post-selection success probability, and the amplitude-encoding fidelity, for both
block encodings; Fig.~\ref{fig:carpet} shows the classical space--time evolution used as that
reference (the corresponding KdV figures are qualitatively identical and included in the
repository). Table~\ref{tab:carleman-resources} gives the exact, extrapolated ($\kappa=8$,
561 phase angles) full-circuit resource cost of \emph{one} QSVT matrix inversion for the systems of
Table~\ref{tab:carleman}, transpiled to the $\{\mathrm{cx},u_3\}$ basis: this is the number that must
be paid at \emph{every} timestep (three timesteps total in each run above), and it is what makes the
Pauli-LCU encoding, despite acting on a matrix that is sparse in the computational basis, the more
expensive of the two by two to three orders of magnitude for this specific matrix family --- because
the Carleman matrix's Kronecker-shift structure gives it a \emph{dense} decomposition in the Pauli
basis (Sec.~\ref{sec:scaling}).

\begin{table}[t]
\caption{Exact extrapolated full-QSVT-circuit resource cost, one matrix inversion, $\kappa=8$
(561 phase angles).}
\label{tab:carleman-resources}
\begin{ruledtabular}
\begin{tabular}{llrrr}
PDE & encoding & qubits & gates & CX gates \\
\hline
Burgers & dense & 6 & 2{,}612{,}767 & 996{,}697 \\
Burgers & Pauli-LCU & 14 & 1{,}716{,}882{,}557 & 803{,}286{,}678 \\
KdV & dense & 6 & 2{,}610{,}531 & 996{,}697 \\
KdV & Pauli-LCU & 13 & 396{,}632{,}217 & 185{,}505{,}488 \\
\end{tabular}
\end{ruledtabular}
\end{table}

\begin{figure*}[t]
\includegraphics[width=\textwidth]{fig_burgers_carleman.pdf}
\caption{Burgers wave propagation, Carleman--QSVT: classical (dashed) vs.\ dense QSVT (solid) vs.\
Pauli-LCU QSVT (dotted).}
\label{fig:carleman-burgers}
\end{figure*}

\begin{figure}[t]
\includegraphics[width=\columnwidth]{fig_burgers_timemarching.pdf}
\caption{Burgers: RMSE, QSVT success probability, and fidelity per timestep, dense vs.\ Pauli-LCU
encoding.}
\label{fig:carleman-time}
\end{figure}

\begin{figure}[t]
\includegraphics[width=\columnwidth]{fig_burgers_carpet.pdf}
\caption{Classical space--time diagram of the Burgers reference solution used above (fine
sub-stepped forward Euler, for visualization only).}
\label{fig:carpet}
\end{figure}

\subsection{What limits the achievable grid size: condition number, not only gate count}
\label{sec:kappa-scaling}
Section~\ref{sec:scaling} shows the \emph{circuit-synthesis} cost of one QSVT layer growing as
$O(N^4)$. A companion analysis in this repository (\texttt{scaling\_analysis/}, following the same
Carleman construction and notation as Ref.~\cite{Setty2026}) shows that a \emph{second}, independent
cost driver grows with grid size as well: the QSVT phase-angle count itself, $O(\kappa\log(1/
\epsilon))$, is set by the block encoding's effective condition number $\kappa_{\rm eff}=\alpha/
\sigma_{\min}$, where $\alpha$ is the block encoding's subnormalization constant --- \emph{not} by
the naive ratio $\sigma_{\max}/\sigma_{\min}$, which that analysis shows tracks $\kappa_{\rm eff}$
poorly. For the Setty-type (Dirichlet) Carleman--Burgers matrix at fixed timestep, $\alpha$ grows with
grid resolution $n_x=1/\Delta x$ through two additive contributions: a diffusive one,
$\lambda_{\rm diff}\propto\nu\Delta t/\Delta x^2\propto n_x^2$, and a convective (CFL-like) one,
$\lambda_{\rm cfl}\propto\Delta t/\Delta x\propto n_x$ multiplied by a combinatorial factor that grows
with the Carleman truncation order $N_T$. In practice the diffusive term dominates at fine
resolution, so refining the grid at fixed $\Delta t$ requires a proportionally larger $\kappa$ budget
(hence more QSVT phase angles and deeper circuits) purely to keep the polynomial approximation to
$1/x$ valid over the encoded matrix's actual spectrum --- a genuine physical constraint on top of, and
independent from, the raw gate-synthesis cost of Sec.~\ref{sec:scaling}. This is exactly the
trade-off we managed by keeping $N_T=2$ and choosing $\nu,\Delta t$ so that the measured
$\kappa_{\rm eff}$ of Table~\ref{tab:carleman} sits inside the $\kappa=8$ budget used throughout.

\section{Method B: Variational natural-gradient time propagation}
\label{sec:variational}

\subsection{Amplitude encoding and ansatz}
The field is represented as $\bm u=\mathrm{amp}\cdot\bm\psi(\bm\theta)$, where $\bm\psi(\bm\theta)$ is
the (always real, unit-norm) statevector of an $n$-qubit \texttt{RealAmplitudes} ansatz with $n$
repetitions of $R_y$ rotations and linear CNOT entanglement, and $\mathrm{amp}\in\R$ is tracked as a
classical scalar. This ansatz has exactly $P=n(n+1)$ free parameters (confirmed for $n=2,\dots,6$ in
Sec.~\ref{sec:scaling}); the initial state is fit once via $\bm\theta_0=\arg\min_{\bm\theta}
\norm{\bm\psi(\bm\theta)-\bm u_0/\norm{\bm u_0}}^2$ with the BFGS optimizer, ten random restarts.

\subsection{The McLachlan variational principle as a discrete-target fit}
Rather than integrating the McLachlan equations of motion for an infinitesimal generator
\cite{McLachlan1964,Li2017,Yuan2019}, we apply the same variational idea to a single, explicit
forward-Euler target: each timestep minimizes
\begin{align}
C(\bm\theta,\mathrm{amp}) &= \norm{\mathrm{amp}\cdot\bm\psi(\bm\theta)-\bm u_{\rm target}}_2^2,\\
\bm u_{\rm target}&=\bm u_{\rm old}+\Delta t\,f(\bm u_{\rm rhs}),
\label{eq:cost}
\end{align}
where $\bm u_{\rm rhs}=\bm u_{\rm old}$ for forward Euler (Burgers) or a separately-solved implicit
state (KdV, Sec.~\ref{sec:implicit-midpoint}). Writing $C_0(\bm\theta)=\braket{\bm\psi(\bm\theta)}
{\bm u_{\rm target}}$, minimizing over $\mathrm{amp}$ at fixed $\bm\theta$ gives
$\mathrm{amp}^\star=C_0(\bm\theta)$ (from $\partial C/\partial\mathrm{amp}=0$); the parameters are
then updated by \emph{natural} gradient descent,
\begin{equation}
\bm\theta\leftarrow\bm\theta-\eta\,M(\bm\theta)^{+}\nabla_{\bm\theta}\big[-2\,\mathrm{amp}^\star\,
C_0(\bm\theta)\big],
\label{eq:ngd}
\end{equation}
where $M^+$ is the pseudo-inverse of the Fubini--Study metric and $\eta$ is chosen by a backtracking
line search. This is the standard cost-function-based variant of the McLachlan principle used
throughout the variational-quantum-simulation literature \cite{Yuan2019}; what is specific to this
work is that every ingredient of Eq.~\eqref{eq:ngd} is measured by a shallow circuit, detailed next.

\subsection{Circuit-level measurement}
\label{sec:circuit-measurement}
Figure~\ref{fig:circuits} shows the general form of the four circuit primitives used throughout this
section (one ansatz repetition shown for $n=3$; the shift gate is drawn for general $n$ by extending
the same staircase pattern).

\begin{figure}[t]
\centering
(a)\\
\begin{quantikz}[column sep=4pt,row sep=4pt]
& \gate{R_y(\theta_1)} & \ctrl{1} & \qw & \qw \\
& \gate{R_y(\theta_2)} & \targ{} & \ctrl{1} & \qw \\
& \gate{R_y(\theta_3)} & \qw & \targ{} & \qw
\end{quantikz}\\[6pt]
(b)\\
\begin{quantikz}[column sep=4pt,row sep=4pt]
\lstick{$q_0$} & \ctrl{2} & \ctrl{1} & \gate{X} & \qw \\
\lstick{$q_1$} & \ctrl{1} & \targ{} & \qw & \qw \\
\lstick{$q_2$} & \targ{} & \qw & \qw & \qw
\end{quantikz}\\[6pt]
(c)\\
\begin{quantikz}[column sep=3pt,row sep=4pt]
\lstick{anc:$\ket0$} & \gate{H} & \gate{X} & \ctrl{1} & \gate{X} & \qw & \ctrl{1} & \gate{H} & \meter{}\\
\lstick{sys} & \qw & \gate{U_{\rm bef}} & \gate{Y_k} & \qw & \gate{U_{\rm mid}} & \gate{Y_l} & \gate{U_{\rm aft}} & \qw
\end{quantikz}\\[6pt]
(d)\\
\begin{quantikz}[column sep=4pt,row sep=6pt]
\lstick{$A$} & \qwbundle{n} & \gate{U(\theta)} & \gate[2]{\text{CX}^{\otimes n}_{A\to B}} & \gate{U(\alpha)^\dagger} & \qw \\
\lstick{$B$} & \qwbundle{n} & \gate{U(\theta)\,S} & & \qw & \qw
\end{quantikz}
\caption{General circuit primitives: (a) one repetition of the $R_y$/CNOT ansatz layer, applied $n$
times in total; (b) the ancilla-free cyclic shift gate $S$ ($n=3$; the staircase of
$(n{-}1{-}k)$-controlled $X$'s targeting qubit $k$, $k=n{-}1,\dots,1$, then $X$ on qubit $0$,
extends to any $n$); (c) the single-ancilla metric off-diagonal circuit, generator insertions $Y_k,Y_l$
between sub-unitaries $U_{\rm bef},U_{\rm mid},U_{\rm aft}$ of the ansatz; (d) the two-register
nonlinear cross-term circuit, $n$-qubit registers $A,B$, no controlled-ansatz gates anywhere.}
\label{fig:circuits}
\end{figure}

\textbf{Metric.} For a single-Pauli-rotation parameter, $M_{kk}=1/4$ exactly, with no circuit needed.
Off-diagonal entries $M_{kl}$ are obtained from a single-ancilla circuit that inserts the two
generators at gate positions $k,l$ and reads $\langle X\rangle$ on the ancilla; we verified this against a
finite-difference evaluation of the exact Fubini--Study metric to $\sim10^{-10}$.

\textbf{Magnitude of each cross term.} $\bm u_{\rm target}$ is a linear combination of overlaps
between $\bm\psi(\bm\theta)$ and $\bm u_{\rm old}$ (identity, each shift term of $F_1$, and the two
building blocks $\chi^{\pm}_i=(S^{\pm1}\bm\psi)_i\,\bm\psi_i$ and
$\mathrm{sq}^{\pm}_i=\big((S^{\pm1}\bm\psi)_i\big)^2$ needed by the skew closure
Eq.~\eqref{eq:f2skew}, $S$ the cyclic shift permutation implemented ancilla-free via a
ripple-carry Toffoli chain. Each is measured by a compute-uncompute circuit whose all-zero
outcome probability equals the overlap's magnitude squared, $P(0)=g^2$ --- for the linear/identity
terms, an ansatz applied forward then a second ansatz applied backward with the shift gate inserted
between them; for the nonlinear terms, two ansatz copies with a shift and a CNOT-based correlation
step, requiring no controlled-ansatz gates at all.

\textbf{The sign problem.} $P(0)$ only gives $|g|$, not its sign, and the standard parameter-shift
derivative of $g$,
\begin{equation}
\frac{\partial g}{\partial\theta_l}=\frac{P(0)\big(\theta_l+\tfrac\pi2\big)-P(0)\big(\theta_l-\tfrac\pi2\big)}{4g},
\label{eq:paramshift}
\end{equation}
is genuinely singular at $g=0$: since $P(0)=g^2$ has a double root there, the numerator is \emph{always}
$O(g^2)\to0$ regardless of the true derivative, so naively treating this as ``derivative $=0$''
silently discards real gradient information (we measured this happening on $\sim8\%$ of evaluations
in a real optimization run --- not a rare edge case). Below a small threshold on $|g|$ we instead take
a classical finite difference of the exact overlap.

\begin{proposition}[exact single-parameter form]
For fixed reference states and a single rotation angle $\theta_l$ (all other parameters held fixed),
$g(\theta_l)=A\cos(\theta_l/2)+B\sin(\theta_l/2)$ exactly, for constants $A,B$ independent of
$\theta_l$.
\end{proposition}
\begin{proof}
The gate acting on $\theta_l$ is $R_y(\theta_l)=\cos(\theta_l/2)I+\sin(\theta_l/2)K$ for a fixed
real matrix $K$ (the generator, embedded via identity on the other qubits). Writing the rest of the
ansatz as fixed unitaries $U_{\rm before},U_{\rm after}$ (independent of $\theta_l$), the state is
$\bm\psi(\theta_l)=U_{\rm after}\big[\cos(\theta_l/2)I+\sin(\theta_l/2)K\big]U_{\rm before}\ket{0}
=\cos(\theta_l/2)\,\bm a+\sin(\theta_l/2)\,\bm b$ for fixed vectors $\bm a,\bm b$. Contracting with a
fixed reference vector $\bm v$ gives $g(\theta_l)=\braket{v}{\psi(\theta_l)}
=\cos(\theta_l/2)\braket va+\sin(\theta_l/2)\braket vb$, of the claimed form with
$A=\braket va,B=\braket vb$.
\end{proof}
This has two consequences we verified numerically to $10^{-10}$: (i) $|\partial g/\partial\theta_l|
\le\tfrac12\norm{\bm v}$, since $A^2+B^2\le\norm{\bm v}^2\norm{\bm\psi}^2=\norm{\bm v}^2$ and the
derivative of the above is bounded by $\tfrac12\sqrt{A^2+B^2}$; and (ii) since every reference vector
used here (an identity/shift-transported ansatz state, or one of the nonlinear vectors
$\chi^\pm,\mathrm{sq}^\pm$) has $\norm{\bm v}_2\le1$, the bound $|\partial g/\partial\theta_l|\le1/2$
holds uniformly for every term.

\subsection{Sign tracking with a rigorous certificate}
\label{sec:signtracker}
Two heuristic attempts at avoiding a classical $O(2^n)$ sign lookup at every evaluation ---
recomputing only when the \emph{current} magnitude looked small, or freezing the sign across a
line-search's candidate points --- both silently corrupted the KdV optimization (9--19\% relative
error, up from $\sim10^{-8}$), because a term's true value can cross zero between two evaluations
without ever being observed near zero at a check point. We instead use the uniform bound just proved:
\begin{proposition}
If $|g(\bm\theta_{\rm last})|>\tfrac12\norm{\bm\theta_{\rm new}-\bm\theta_{\rm last}}_1$, then $g$
cannot vanish anywhere on the straight line between $\bm\theta_{\rm last}$ and $\bm\theta_{\rm new}$.
\end{proposition}
\begin{proof}
Along the line $\bm\theta(t)=\bm\theta_{\rm last}+t(\bm\theta_{\rm new}-\bm\theta_{\rm last})$,
$t\in[0,1]$, the chain rule and the uniform bound give
$\big|\tfrac{d}{dt}g(\bm\theta(t))\big|\le\sum_l\big|\tfrac{\partial g}{\partial\theta_l}\big|\,
|\theta_{{\rm new},l}-\theta_{{\rm last},l}|\le\tfrac12\norm{\bm\theta_{\rm new}-\bm\theta_{\rm
last}}_1$. Integrating, $|g(\bm\theta(t))-g(\bm\theta_{\rm last})|\le t\cdot\tfrac12\norm{\bm\theta_{\rm
new}-\bm\theta_{\rm last}}_1$ for all $t\in[0,1]$, so if the right-hand side (at $t=1$) is smaller
than $|g(\bm\theta_{\rm last})|$, then $g(\bm\theta(t))$ cannot reach $0$ for any $t\in[0,1]$.
\end{proof}
This is a \emph{proven} certificate, not a heuristic: a term's sign is carried forward at zero extra
circuit cost whenever it holds, and falls back to one classical evaluation only when it fails.
Measured on a real Burgers run, this reduced classical calls from one full $O(2^n)$ lookup per trial
point (thousands per timestep) to $\sim55$ per timestep, with the resulting optimization matching the
always-classical baseline to full numerical precision on both PDEs and multiple random seeds.

\subsection{Implicit midpoint for KdV}
\label{sec:implicit-midpoint}
Section~\ref{sec:forward-euler-injects-energy} proved forward Euler injects energy into KdV on every
step. The implicit midpoint rule, $\bm u_{n+1}=\bm u_n+\Delta t\,f\!\big(\tfrac{\bm u_n+\bm u_{n+1}}2\big)$,
does not have this defect:
\begin{proposition}
Implicit midpoint conserves $\norm{\bm u}_2^2$ exactly, for any $\Delta t$, whenever
$\bm v^{\!\top}f(\bm v)=0$ for all $\bm v$.
\end{proposition}
\begin{proof}
Let $\bm u_{\rm mid}=\tfrac12(\bm u_n+\bm u_{n+1})$, so $\bm u_{n+1}-\bm u_n=\Delta t\,f(\bm u_{\rm
mid})$. Then $\norm{\bm u_{n+1}}_2^2-\norm{\bm u_n}_2^2=(\bm u_{n+1}-\bm u_n)^{\!\top}(\bm u_{n+1}
+\bm u_n)=\Delta t\,f(\bm u_{\rm mid})^{\!\top}\cdot2\bm u_{\rm mid}=2\Delta t\,\bm u_{\rm mid}^{\!\top}
f(\bm u_{\rm mid})=0$, using the hypothesis applied to $\bm u_{\rm mid}$ itself.
\end{proof}
We solve the (small, classical, $N$-dimensional) nonlinear equation for $\bm u_{n+1}$ once per
timestep, fit a second ansatz encoding to $\bm u_{\rm mid}$ by the same BFGS procedure as the initial
state, and use it as $\bm u_{\rm rhs}$ in Eq.~\eqref{eq:cost} while $\bm u_{\rm old}$ still supplies
the identity term. Confirmed to $10^{-10}$ classically ($\norm{\bm u(t)}_2=1.0000000000$ at every one
of three steps, versus $1.00\to1.19\to1.42\to1.72$ under forward Euler), and to $\sim10^{-7}$ for the
quantum-fitted state.

\subsection{Verification and results}
Table~\ref{tab:variational} summarizes both systems, and Table~\ref{tab:variational-resources} gives
the exact transpiled cost of each circuit type in Fig.~\ref{fig:circuits}, at the qubit count actually
used by each system. Figures~\ref{fig:burgers-var} and~\ref{fig:kdv-var} show the resulting time
propagation: the ideal (statevector) quantum curve is visually indistinguishable from the classical
target at every panel, and the noisy-hardware curve (FakeMarrakesh/FakeSherbrooke calibration
snapshot, shot-based readout with device noise, sign taken from the ideal simulation since only
$\sqrt{P(0)}$ is measurable this way) degrades gracefully.

\begin{table}[t]
\caption{Variational natural-gradient systems solved ($\Delta t=0.1$, $n_{\rm iter}=150$).}
\label{tab:variational}
\begin{ruledtabular}
\begin{tabular}{lccc}
PDE & $n$ & $P$ & rel.\ $L^2$ err.\ \\
& (qubits) & & ideal / noisy \\
\hline
Burgers ($\nu=0.01$) & 3 & 12 & $2\text{--}4\!\times\!10^{-6}$ / $0.17\text{--}0.20$ \\
KdV ($\delta=0.0625$) & 2 & 6 & $2\!\times\!10^{-8}\text{--}2\!\times\!10^{-7}$ / $0.07\text{--}0.13$ \\
\end{tabular}
\end{ruledtabular}
\end{table}

\begin{table}[t]
\caption{Exact transpiled cost of each circuit primitive (Fig.~\ref{fig:circuits}), one instance,
$\{\mathrm{cx},u_3\}$ basis.}
\label{tab:variational-resources}
\begin{ruledtabular}
\begin{tabular}{llrrr}
PDE & circuit & qubits & gates & depth \\
\hline
Burgers ($n=3$) & identity/linear & 3 & 128 & 68 \\
Burgers ($n=3$) & nonlinear & 6 & 351 & 148 \\
Burgers ($n=3$) & metric & 4 & 85 & 40 \\
KdV ($n=2$) & identity/linear & 2 & 29 & 17 \\
KdV ($n=2$) & nonlinear & 4 & 140 & 70 \\
KdV ($n=2$) & metric & 3 & 55 & 33 \\
\end{tabular}
\end{ruledtabular}
\end{table}

\begin{figure*}[t]
\includegraphics[width=\textwidth]{fig_burgers_variational.pdf}
\caption{Burgers wave propagation, classical (dashed) vs.\ ideal quantum (solid, matches to
$\sim10^{-6}$) vs.\ noisy hardware readout (dotted).}
\label{fig:burgers-var}
\end{figure*}

\begin{figure*}[t]
\includegraphics[width=\textwidth]{fig_kdv_variational.pdf}
\caption{KdV wave propagation with the implicit-midpoint target of
Sec.~\ref{sec:implicit-midpoint}; $\norm{u}_2$ is conserved to $10^{-10}$ classically and matched to
$\sim10^{-7}$ by the quantum fit.}
\label{fig:kdv-var}
\end{figure*}

\section{Resource scaling}
\label{sec:scaling}

We measured, directly from our own implementation (not estimated), how each pathway's resources
scale with problem size, transpiling every circuit against the \texttt{FakeSherbrooke} calibration
snapshot at optimization level 3.

\textbf{Variational.} At $n=2,\dots,6$ qubits ($N=2^n=4,\dots,64$), every measured quantity is
polynomial in $n$, hence \emph{polylogarithmic} in $N$: the ansatz parameter count is the exact
quadratic $P=n(n+1)=O(\log^2N)$ (confirmed exactly at every $n$ tested), transpiled gate count and
depth of all three circuit types fit power laws with exponents in $[1.2,2.6]$ (so
$O(\mathrm{poly}(\log N))$), and qubit count is exactly $O(n)=O(\log N)$ ($n$ for the identity/linear
circuit, $2n$ for the nonlinear one, $n+1$ for the metric) --- Table~\ref{tab:bigo} summarizes.

\textbf{Carleman--QSVT, dense encoding.} The dense block encoding's per-layer gate count fits
$\propto2^{2.03\,n_w}=O(4^{n_w})$ against its own qubit count $n_w$ --- the well-known generic cost of
synthesizing an arbitrary $n_w$-qubit unitary with no exploitable structure \cite{Shende2006}. Since
$n_w=O(\log(N+N^2))=O(\log N)$ at our truncation order $N_T=2$, this is $O(N^{4})$ in grid size $N$:
polynomial, but far steeper than the variational method's polylogarithmic scaling, and it must
additionally be multiplied by the QSVT phase-angle count and by the number of timesteps.

\textbf{Carleman--QSVT, Pauli-LCU encoding.} The number of nonzero Pauli terms needed grows noisily
but substantially faster than the matrix dimension itself, because the Kronecker-shift structure of
Eq.~\eqref{eq:carleman} is sparse in the computational basis but dense in the Pauli basis --- which
is why we measured the ostensibly-sparse LCU encoding to require \emph{more} gates than the dense one
at matched $N$ (Table~\ref{tab:carleman-resources}).

\begin{table*}[t]
\caption{Measured asymptotic scaling, this implementation ($n$=qubits, $N=2^n$).}
\label{tab:bigo}
\begin{ruledtabular}
\begin{tabular}{lll}
Pathway & Quantity & Scaling \\
\hline
Variational & ansatz parameters $P$ & $\Theta(n^2)=\Theta(\log^2N)$ \\
Variational & qubits (any circuit) & $O(n)=O(\log N)$ \\
Variational & gates, depth & $O(n^{2.6})=O(\mathrm{poly}(\log N))$ \\
Carleman, dense encoding & gates per QSVT layer & $O(4^{n_w})=O(N^4)$ \\
Carleman, Pauli-LCU encoding & nonzero Pauli terms & $\omega(N)$, super-linear \\
\end{tabular}
\end{ruledtabular}
\end{table*}

\textbf{A different bottleneck under a sparse (rather than dense) block encoding.} The $O(N^4)$ above
is specific to our \emph{dense}-unitary-synthesis encoding. Using instead a sparse arithmetic oracle
(indexing the $O(N_T)$ nonzero entries per row directly, rather than synthesizing a dense unitary),
the companion analysis in \texttt{scaling\_analysis/} derives the alternative asymptotic sequential
two-qubit depth, for the Dirichlet-BC construction of Ref.~\cite{Setty2026} under standard QSVT
assumptions,
\begin{equation}
D_{\rm QSVT}=\widetilde{O}\!\left(\frac{\kappa_V\,T\,N_T^{3}}{\mathrm{Re}\,\epsilon}\right),
\label{eq:sparse-depth}
\end{equation}
with $\kappa_V$ the Carleman matrix's eigenvector condition number and $\epsilon$ the target error ---
here the bottleneck is the QSVT polynomial degree (set by the condition number, which worsens as
$\mathrm{Re}\to0^+$ towards the edge of Carleman convergence) rather than gate-synthesis cost, and it
grows with $\mathrm{Re}^{-1}$, not with $N$ directly. That analysis's worked instance at the largest
Reynolds number ($\mathrm{Re}=1$) for which convergence is provable found $D_{\rm QSVT}\sim10^{12}$
two-qubit gates against $\sim10^6$ classical floating-point operations for the same accuracy on the
same instance --- a reminder that, at the modest Reynolds numbers a first-order Carleman truncation
can provably handle, the classical solve remains many orders of magnitude cheaper, independent of
which block-encoding strategy is used.

\section{Discussion}
\label{sec:discussion}

The two pathways trade off in the expected direction: Carleman--QSVT solves a genuinely linear
problem exactly (up to truncation and QSVT polynomial-approximation error) but pays an $O(N^4)$
circuit cost that made even a modest grid ($N=8$) impractical to fully synthesize within our
compute budget; the variational method's circuits scale polylogarithmically in $N$ but solve a
nonconvex optimization problem afresh at every timestep, with no convergence guarantee. Because of
this we did not force the two pathways onto identical grids --- Carleman uses the largest $N$ we
could fully synthesize and verify, and the variational method uses the smallest $N$ needed to make
its own coarse-grid physics honest (Sec.~\ref{sec:closure}'s closure only guarantees the required
$L^2$ bounds; it does not by itself guarantee that a coarse grid resolves the solution's shape well).

One gap remains open rather than closed: Sec.~\ref{sec:signtracker}'s certificate reduces, but does
not eliminate, the classical $O(2^n)$ sign lookup in the variational method. We built and verified
(to $10^{-16}$) an alternative that removes it entirely, by measuring the \emph{signed} overlap
directly with a Hadamard test (ancilla-controlling every gate in each compute-uncompute circuit); we
did not adopt it because doing so increases transpiled circuit depth roughly ninefold to elevenfold
for the same circuit count, and the certificate of Sec.~\ref{sec:signtracker} already removes the
classical dependency from all but a small, measured fraction of evaluations. Making that residual
fraction provably $O(\mathrm{poly}(n))$ rather than merely empirically small remains open.

\section{Conclusion}

We presented and fully verified two independent quantum pathways for the viscous Burgers and KdV
equations on a periodic domain, sharing a single discretization whose convective closure we proved,
in three lines, to be exactly energy-conserving --- the property responsible for both PDEs meeting
their required $L^2$ norm bounds. We proved that naively time-stepping KdV's conservative ODE with
forward Euler injects energy on every step, proved that implicit midpoint does not, and implemented
the fix. We proved a Lipschitz bound on the variational method's parameter dependence and used it to
build a certified (not heuristic) alternative to full classical sign recomputation, verified against
two independent, previously-broken heuristic attempts. Finally, we measured, directly from our own
circuits rather than by estimation, that the variational method's resources scale
polylogarithmically in grid size while the Carleman--QSVT dense block encoding scales as $O(N^4)$,
quantifying precisely the trade-off between the two approaches.

\begin{acknowledgments}
We thank the organizers of the 2026 Niels Bohr Quantum Summer School Challenge on Quantum Simulation
of Nonlinear Partial Differential Equations, Center for Quantum Mathematics, University of Southern
Denmark.
\end{acknowledgments}

\section*{AI disclosure}
Substantial parts of this work's implementation, debugging, numerical verification, and this
manuscript were produced with the assistance of Claude (Anthropic), used interactively under the
direction and review of the authors. All physical claims, discretization choices, and numerical
results in this paper were independently verified against classical reference calculations rather
than accepted on the model's assertion alone; every proof was checked step by step; every cited
reference was checked against its original source before inclusion.

\bibliographystyle{apsrev4-2}
\begin{thebibliography}{99}

\bibitem{Setty2026} A.~Setty, ``A quantum linear systems pathway for solving differential equations,''
J. Phys. A: Math. Theor. \textbf{59}, 185303 (2026).

\bibitem{Liu2021} J.-P.~Liu, H.~{\O}.~Kolden, H.~K.~Krovi, N.~F.~Loureiro, K.~Trivisa, A.~M.~Childs,
``Efficient quantum algorithm for dissipative nonlinear differential equations,'' Proc. Natl. Acad.
Sci. U.S.A. \textbf{118}, e2026805118 (2021).

\bibitem{Gilyen2019} A.~Gily{\'e}n, Y.~Su, G.~H.~Low, N.~Wiebe, ``Quantum singular value
transformation and beyond: exponential improvements for quantum matrix arithmetics,'' in
\emph{Proceedings of the 51st Annual ACM SIGACT Symposium on Theory of Computing} (STOC 2019),
pp.~193--204.

\bibitem{McLachlan1964} A.~C.~McLachlan, ``A variational solution of the time-dependent Schr\"odinger
equation,'' Mol. Phys. \textbf{8}, 39 (1964).

\bibitem{Li2017} Y.~Li, S.~C.~Benjamin, ``Efficient Variational Quantum Simulator Incorporating
Active Error Minimization,'' Phys. Rev. X \textbf{7}, 021050 (2017).

\bibitem{Yuan2019} X.~Yuan, S.~Endo, Q.~Zhao, Y.~Li, S.~C.~Benjamin, ``Theory of variational quantum
simulation,'' Quantum \textbf{3}, 191 (2019).

\bibitem{Mitarai2018} K.~Mitarai, M.~Negoro, M.~Kitagawa, K.~Fujii, ``Quantum circuit learning,''
Phys. Rev. A \textbf{98}, 032309 (2018).

\bibitem{Schuld2019} M.~Schuld, V.~Bergholm, C.~Gogolin, J.~Izaac, N.~Killoran, ``Evaluating analytic
gradients on quantum hardware,'' Phys. Rev. A \textbf{99}, 032331 (2019).

\bibitem{Costa2025} P.~C.~S.~Costa, P.~Schleich, M.~E.~S.~Morales, D.~W.~Berry, ``Further improving
quantum algorithms for nonlinear differential equations via higher-order methods and rescaling,''
npj Quantum Inf. \textbf{11}, 141 (2025).

\bibitem{Harten1983} A.~Harten, ``High resolution schemes for hyperbolic conservation laws,''
J. Comput. Phys. \textbf{49}, 357 (1983).

\bibitem{Tadmor1987} E.~Tadmor, ``The numerical viscosity of entropy stable schemes for systems of
conservation laws. I,'' Math. Comp. \textbf{49}, 91 (1987).

\bibitem{Shende2006} V.~V.~Shende, S.~S.~Bullock, I.~L.~Markov, ``Synthesis of quantum-logic
circuits,'' IEEE Trans. Comput.-Aided Des. Integr. Circuits Syst. \textbf{25}, 1000 (2006).

\bibitem{Fornberg1988} B.~Fornberg, ``Generation of finite difference formulas on arbitrarily spaced
grids,'' Math. Comp. \textbf{51}, 699 (1988).

\bibitem{LeVeque2007} R.~J.~LeVeque, \emph{Finite Difference Methods for Ordinary and Partial
Differential Equations} (SIAM, Philadelphia, 2007).

\bibitem{LeVeque2002} R.~J.~LeVeque, \emph{Finite Volume Methods for Hyperbolic Problems} (Cambridge
University Press, Cambridge, 2002).

\bibitem{DDCL2026} Niels Bohr Quantum Summer School 2026, ``Analog Simulation of Burgers Dynamics in a
Digital Degenerate Cavity Laser,'' Center for Quantum Mathematics, University of Southern Denmark.

\end{thebibliography}

\end{document}
